%==========================================================================
\section{Related Work}\label{sec:related}

\subsection{FPGA statevector simulators and emulators}\label{subsec:related-fpga}

Early FPGA quantum-circuit emulators held the whole statevector in on-chip memory and were limited to roughly ten to sixteen qubits~\cite{khalid2004fpga,aminian2008fpga,lee2016fpga,pilch2019fpga}. Mahmud et al.\ scaled reconfigurable emulation to higher precision and throughput on a host-attached FPGA platform~\cite{mahmud2019scaling}. More recent work targets data-centre cards with high-bandwidth memory: Moawad et al.\ studied FPGA acceleration of full-state simulation for computational-fluid-dynamics circuits~\cite{moawad2022cfd} and proposed iteration-scheduling memory access patterns that reduce the work of controlled gates~\cite{moawad2024iteration}; Bennakhi et al.\ presented an HLS state-vector simulator (Q2SV)~\cite{bennakhi2025q2sv}; Tran et al.\ reported a resource-efficient state-vector accelerator on an Alveo U280 (QEA)~\cite{tran2025qea} and, with HPQEA, an HBM-based emulator on the same card that reaches 30 qubits and is compared with an NVIDIA A100~\cite{tran2025hpqea}; and Suzuki et al.\ combined CPU and FPGA in a quantum machine-learning simulator~\cite{suzuki2023hybrid}. Cicero et al.\ survey simulation methods and their acceleration on CPUs, GPUs and FPGAs~\cite{cicero2024review}. Thirty qubits on one HBM card are therefore not a record, and we make no such claim. Our work differs in three respects: (i) a measured design-space study of the HBM memory system itself (pseudo-channels, engines, port width, layout, exchange), (ii) a single datapath primitive that covers every one- and two-qubit gate and multi-controlled gates natively, and (iii) a measured cross-platform comparison with a published, reproducible timing boundary.

\subsection{HBM on FPGAs}\label{subsec:related-hbm}

Wang et al.\ benchmarked the HBM subsystem of the Alveo U280 and showed that the achieved bandwidth (up to 425~GB/s) depends strongly on access pattern and on how the pseudo-channels are used~\cite{wang2020shuhai}. Choi et al.\ showed that HLS designs reach only a fraction of the HBM bandwidth unless burst length, outstanding requests and channel assignment are tuned~\cite{choi2020hlshbm}, and proposed a customised HLS interconnect for designs that need all-to-all channel access~\cite{choi2021hbmconnect}. Section~\ref{subsec:res-ceilings} reports the corresponding measurements on the U55C used here; they set the ceiling against which every simulator result is compared.

\subsection{GPU and CPU statevector simulators}\label{subsec:related-gpu}

GPU statevector simulation is the established approach for 20 to 34 qubits on a single device: Qiskit Aer~\cite{qiskit2024}, NVIDIA's cuStateVec within the cuQuantum SDK~\cite{bayraktar2023cuquantum} and CUDA-Q~\cite{nvidia2024cudaq} are widely used. On CPUs, QuEST~\cite{jones2019quest} and Qulacs~\cite{suzuki2021qulacs} provide multi-threaded statevector engines. Distributed statevector simulation trades memory capacity for inter-node communication~\cite{chen2018quantum,pednault2017breaking}; tensor-network methods~\cite{vidal2003efficient,markov2008simulating}, used for example by the qFlex simulator~\cite{villalonga2019flexible}, and stabiliser~\cite{gottesman1998heisenberg,aaronson2004improved} methods avoid the exponential state for restricted circuit classes.
